%!TEX program = xelatex
\documentclass[UTF8,a4paper,zihao=-4,fontset=fandol]{ctexart}
\usepackage[margin=24mm,headheight=15pt]{geometry}
\usepackage{amsmath,amssymb,amsthm,fancyhdr,xurl,needspace}
\usepackage[backend=biber,style=numeric,sorting=none,maxbibnames=99]{biblatex}
\addbibresource{../bibliography/preview_references.bib}
\theoremstyle{remark}
\newtheorem*{remark}{注}
\setlength{\parskip}{4pt}
\emergencystretch=2em
\allowdisplaybreaks[3]
\pagestyle{fancy}\fancyhf{}
\fancyhead[L]{社区感染阈值与医疗容量的条件性联系}
\fancyhead[R]{修改阅读稿}
\fancyfoot[L]{局部替换内容；不是论文全文}
\fancyfoot[R]{\thepage}
\makeatletter
\newlabel{sec:model:strategy}{{2.2}{}}
\newlabel{eq:model:full}{{2.1}{}}
\newlabel{sec:joint:compare-numerics}{{7.2}{}}
\newlabel{sec:xian:threshold-scale}{{8.3}{}}
\makeatother
\begin{document}
\begin{center}
{\Large 医疗容量条件性联系：修改内容阅读稿}\par
\vspace{5pt}2026年10月5日\par
\end{center}
本文件按正文位置展示拟替换段落，未覆盖整篇论文。段落中指向原稿的章节与公式引用沿用其原位置；本阅读稿自身的公式编号与参考文献编号不作为论文最终编号。全部计算结果沿用原稿，本轮没有新增医疗需求模拟或临床标定。
\section*{摘要开头与引言衔接}
避免医疗需求过载是限制社区感染高峰的重要动机，而实现同一感染上限的不同措施组合可能产生不同代价。本文在无隔离易感者回流的 SIQR 型模型中，通过全部新增感染流量的上界，给出共同医疗需求假设下平均 ICU 需求不超过容量的充分条件。在此基础上，以给定且满足实施条件的非隔离感染人数上限为控制目标，研究首次达到阈值后保持感染人数不变、并在常规传播临界状态退出的一类控制。

\medskip
在区分隔离状态的模型中，选择非隔离感染者作为控制对象有直接的传播学依据。Zhou 等\cite{Zhou2024SPCI}指出，隔离区外的传播主要由尚未被隔离或追踪的感染者造成，其中包括不易被发现的无症状感染者。在本文模型中，进入非隔离感染仓室 $I$ 和隔离感染仓室 $I_q$ 的全部新增感染均由社区中的 $I$ 引起。因此，在接触率有界的条件下，限制 $I$ 不仅控制社区感染峰值，也为两类病例的总新增感染流量提供统一上界。隔离感染者虽不再参与模型中的传播，仍可能需要救治，不能将其医疗需求排除在外。Zhang 等\cite{Zhang2026Behavior}在含住院过程的行为--传播模型中，也使用感染现患率上界讨论医疗压力显现之前的干预。

本文以避免医疗需求过载为现实动机，以社区非隔离感染峰值为直接控制对象。第\ref{sec:model:strategy}节通过全部新增感染流量，说明在共同的医疗需求关系和初始负担条件下，怎样选取足够低的社区感染阈值，使平均 ICU 需求不超过可用容量。这一联系不增加传播仓室，也不将恒定社区感染人数等同于 ICU 满负荷运行。后续分析以给定、可实施的阈值为基础，研究规定三阶段策略的解析结构和两种措施的分配；具体地区的临床容量校准与这一控制分析分别处理。
\clearpage
\section*{第2.2节：阈值的条件性医疗解释}
本文以社区非隔离感染人数为直接控制对象，要求全过程满足 $I(t)\le\eta$，其中 $\eta>0$ 为给定上限，并记 $\theta=\eta/N$。Miclo 等\cite{Miclo2022ICU}通过比例关系将感染人数上限与 ICU 容量联系起来；Angulo 等\cite{Angulo2021}则先给定允许现患率，再在应用中依据医疗资源说明其取值。由于本文区分非隔离与隔离感染者，容量解释还需要计入两类病例的共同来源。以下说明社区感染上限与医疗容量的联系，再研究给定阈值下的控制策略。

\begin{remark}
在模型\eqref{eq:model:full}中，两类感染仓室的新增流入之和为
\[
\frac{\beta c(t)(1-q(t))S(t)I(t)}{N}
+\frac{\beta c(t)q(t)S(t)I(t)}{N}
=\frac{\beta c(t)S(t)I(t)}{N}.
\]
若全过程有 $0\le c(t)\le c_0$ 和 $I(t)\le\eta$，则由 $S(t)\le S_0$ 得
\[
0\le\frac{\beta c(t)S(t)I(t)}{N}
\le\frac{\beta c_0S_0}{N}\eta.
\]
这给出全部新增感染流量的统一上界，而不只约束留在社区的病例。

为说明该上界与医疗容量的联系，附加如下医疗需求关系。设 $h(a)\ge0$ 为一例新感染者在感染后第 $a$ 天所需的平均 ICU 床位数；假定该函数可测，且对进入 $I$、进入 $I_q$ 的两类病例、不同感染时点及所比较的干预分配共同适用。设
\[
\int_0^\infty h(a)\,da=pL<\infty,
\]
其中 $p\in(0,1]$ 是一例模型新感染者此后需要 ICU 的概率，$L>0$ 是需要 ICU 者累计占用时长的均值。$h$ 可以包含感染至入院的延迟，不要求指定其具体分布。由既往流入与停留过程计算床位占用的建模方法可参见文献\cite{Baas2021Occupancy}。

记初始病例在此后产生的平均需求为 $D_0(t)$，并假定对所比较策略均有 $0\le D_0(t)\le D_{\mathrm{init}}<\infty$。该项包括初始时刻已经感染但以后才需要 ICU 的病例，不能仅用初始在床人数代替。未受容量拒收限制的平均需求定义为
\[
D(t)=D_0(t)+\int_0^t h(t-u)
\frac{\beta c(u)S(u)I(u)}{N}\,du.
\]
由非负性及前述流量界，直接得到
\[
\begin{aligned}
D(t)
&\le D_{\mathrm{init}}+
\frac{\beta c_0S_0}{N}\eta\int_0^t h(a)\,da\\
&\le D_{\mathrm{init}}+pL\frac{\beta c_0S_0}{N}\eta,
\qquad t\ge0.
\end{aligned}
\]
\Needspace{9\baselineskip}
设 $C>D_{\mathrm{init}}$ 为扣除其他病种占用和安全预留后可用于本病的容量，其中本病初始病例的负担由 $D_{\mathrm{init}}$ 单独计入、不重复扣除。因而
\[
\eta\le\frac{C-D_{\mathrm{init}}}{pL\,\beta c_0S_0/N}
\]
是上述平均需求满足 $D(t)\le C$ 的充分条件。两类病例均已计入总流量，不需要假定隔离感染者没有重症风险。该条件覆盖控制前、控制期和退出后；它不要求需求峰值与感染峰值同时出现，也不表示 $I=\eta$ 时 ICU 恰好满载。这里的结论针对所规定的平均需求关系，不是实际随机占用的逐样本保证；不满足该充分条件也不能据此判定必然超容。
\end{remark}

上述医疗需求关系用于解释阈值的选取，不加入传播方程、人口守恒式或成本目标。本文不估计 $h,p,L$ 或初始病例的医疗需求；应用中的比例参考另行说明，不将其与这里的充分条件等同。以下对满足初值、触发及实施条件的给定 $\eta$ 分析三阶段控制。各策略使用同一阈值，不按各自的感染分流结果重新调整。

\medskip
从医疗容量选择阈值与检验干预能力是两个步骤。若按上述平均需求充分条件选取 $\eta$，还须满足下文的触发条件；有额外能力限制时，再分别检查仅隔离或联合控制的实施判据。满足这些条件后，同一医疗需求上界适用于所比较的全部允许分配，但不意味着它们具有相同的实际需求轨迹。后文只在共同启动状态、阈值及退出目标下比较控制，不把这一充分条件作为新的时变状态约束加入成本最小化。
\clearpage
\section*{第8.3节：设计情景的解释}
本节将传播参数的拟合与感染阈值的设定分开。西安市 2021 年末常住人口为 $N=13{,}163{,}000$\cite{XianStat2021}。现有资料尚不足以确定与本模型全部新增感染口径一致的 ICU 需求函数、初始病例未来负担及同期可用容量，因而本节没有按第\ref{sec:model:strategy}节的充分条件校准阈值。以下保留一个文献比例参考及其附近的设计情景，用于比较传播与干预指标。

为与已有情景保持一致，仅在本小节使用参考系数 $\rho_{\mathrm{ICU}}>0$，按形式关系 $\eta=C/\rho_{\mathrm{ICU}}$、$\theta=(C/N)/\rho_{\mathrm{ICU}}$ 说明阈值量级。陈胤孜等\cite{Chen2021ICUForecast}基于既往资源资料预测，2021 年全国每10万人综合 ICU 床位约为 $4.37$ 张；它不是西安同期可用容量的实测值。Angulo 等\cite[补充材料 S4.1]{Angulo2021}对早期 COVID-19 采用需要重症监护的感染者比例名义值 $0.60\times0.15\times0.28=0.0252$。暂将全国预测量作为人均容量参考、不扣除既有占用，并把该名义比例移用于本节的参考系数，得到
\[
\theta=\frac{4.37/10^5}{0.0252}\approx1.73\times10^{-3},
\qquad \eta\approx2.28\times10^4.
\]
这一计算只给出比例情景的量级。$\rho_{\mathrm{ICU}}$ 没有依据共同需求函数及初始负担估计，不是第\ref{sec:model:strategy}节充分条件中的 $pL\beta c_0S_0/N$，也不是所有策略下实际占用与 $I(t)$ 的固定比值。以确诊病例或全部感染者为分母得到的 ICU 入院风险本身就可能不同\cite{Pung2023Undetected}，上述移植不能作为本地临床校准。

为比较这一量级附近的控制结果，保留代表性设计情景
\begin{equation}
\eta=0.002N=26{,}326.
\end{equation}
这不是上述名义换算值的保守取整。在 $\rho_{\mathrm{ICU}}=0.0252$ 的比例情景中，它对应每10万人 $5.04$ 张、全市约 $663$ 张的名义容量；这些数值不作为西安实际可用资源。第\ref{sec:joint:compare-numerics}节及本节采用该阈值的算例，均为同一社区感染控制目标下的策略比较，尚未核查第\ref{sec:model:strategy}节的容量充分条件。它们既不构成实际 ICU 不超容的证明，也不能据未校准而判定必然超容。阈值变化的影响仍由后续敏感性分析说明。
\clearpage
\section*{第10.2节：医疗联系与适用范围}
本文将避免医疗需求过载作为现实动机，但直接控制的是社区非隔离感染人数。第\ref{sec:model:strategy}节给出的联系经过全部新增感染流量，因此覆盖后来进入 $I_q$ 的病例，而不是把隔离病例排除在医疗需求之外。这与 Zhou 等\cite{Zhou2024SPCI}在完整模型中保留两类病例住院流入的处理相容；该文支持控制非隔离感染峰值的传播学理由，并未替本文建立容量充分条件。Zhang 等\cite[注记1]{Zhang2026Behavior}的感染--住院比例是其模型的平衡态关系，也不作为本文时变轨迹的固定换算。

上述平均需求界依赖共同的需求函数以及初始病例未来需求的上界，采用 $c_0S_0$ 得到对全部允许分配统一的充分条件，可能偏保守。它不要求隔离与非隔离感染人数成固定比例，但假定两类病例及不同分配采用相同的感染后需求关系；若医疗需求随病例来源、感染时点或干预措施改变，则需要重新建立共同上界。满足条件表示所设平均需求不超过容量，不表示真实随机占用的超载概率为零，也不能据此量化死亡率变化。第\ref{sec:xian:threshold-scale}节及联合控制比较中的 $\theta=0.002$ 仍是设计情景，未完成这一容量校准。$I+I_q$ 的峰值只反映总现存感染人数，不直接等同于住院或 ICU 占用。

\medskip
将第\ref{sec:model:strategy}节的充分条件用于具体疫情，仍需与模型新感染口径一致的临床需求资料、初始病例负担和可用容量。若以后把实际时变医疗占用直接作为优化约束，还须检查其历史依赖是否改变原来的逐状态成本分解。
\section*{结论开头}
本文以避免医疗需求过载为现实动机，通过全部新增感染流量给出社区感染上限与平均 ICU 需求之间的条件性联系。在给定、可实施的非隔离感染人数上限下，含接触追踪的 SIQR 型模型的阈值控制过程由常规阶段首次积分和恒定感染条件连接起来；主要分析仍针对这一规定过程的解析结构与干预分配。
\clearpage
\printbibliography[title={阅读稿参考文献}]
\end{document}
